\documentclass[a4paper,10pt,brazil]{article}
\usepackage{provastyle}
\usepackage{menukeys}

\begin{document}
\makeexamtitlepage{UFOP}{BCC222: Programação Funcional}{José Romildo}{2026/1}{Primeira Prova}{08}
\begin{multicols}{2}
\questionTitle{1}{Assinatura de função com dois Int e String}
\begin{flushright}
  \footnotesize
  \textit{\textbf{4: Tipos de Dados}}\\
  \textit{c04-funcao-assinatura-01}\\

\end{flushright}
Qual das seguintes assinaturas de tipo representa uma função que recebe dois \haskellinline{Int} e uma \haskellinline{String} e retorna um \haskellinline{Bool}?

\begin{enumerate}
  \item \haskellinline{(Int, Int, String) -> Bool}

  \item \haskellinline{Int -> Int -> String -> Bool}

  \item \haskellinline{[Int, Int, String] -> Bool}

  \item \haskellinline{Int -> Int -> String, Bool}

  \item \haskellinline{Bool -> Int -> Int -> String}


\end{enumerate}

\questionTitle{2}{Avaliação parametrizada de expressão let}
\begin{flushright}
  \footnotesize
  \textit{\textbf{3: Expressões e definições}}\\
  \textit{c03-let-escopo-01}\\

\end{flushright}
Qual é o valor da expressão abaixo?

\begin{haskellcode}
let n = 3; q = n ^ 2 in q + n
\end{haskellcode}

\begin{enumerate}
  \item \haskellinline|9|.

  \item \haskellinline|3|.

  \item \haskellinline|12|.

  \item A expressão é inválida porque \haskellinline|let| só pode aparecer no início de uma definição de função.

  \item \haskellinline|27|.


\end{enumerate}

\questionTitle{3}{Próxima equação da função}
\begin{flushright}
  \footnotesize
  \textit{\textbf{5: Expressão condicional}}\\
  \textit{c05-guardas-exaustividade-03}\\

\end{flushright}
Considere a definição:
\begin{haskellcode}
h x
  | x > 0 = "Positivo"
h x = "Não positivo"
\end{haskellcode}
Qual é o resultado de \haskellinline|h 0|?

\begin{enumerate}
  \item Ocorre um erro, pois a primeira equação não tem \haskellinline|otherwise|.

  \item \haskellinline|"Não positivo"|

  \item O compilador rejeita a definição por possuir duas equações para a mesma função.

  \item \haskellinline|0|

  \item \haskellinline|"Positivo"|


\end{enumerate}

\questionTitle{4}{Inferência em discriminante de equação quadrática}
\begin{flushright}
  \footnotesize
  \textit{\textbf{7: Polimorfismo}}\\
  \textit{cap07-inferencia-possui-raizes}\\

\end{flushright}
Considere:

\begin{minted}{haskell}
possuiRaizesReais a b c = b^2 - 4*a*c >= 0
\end{minted}

Qual assinatura é a mais geral?

\begin{enumerate}
  \item \mintinline{haskell}{possuiRaizesReais :: Floating a => a -> a -> a -> [a]}
  \item \mintinline{haskell}{possuiRaizesReais :: Num a => a -> a -> a -> Bool}
  \item \mintinline{haskell}{possuiRaizesReais :: Bool -> Bool -> Bool -> Bool}
  \item \mintinline{haskell}{possuiRaizesReais :: (Num a, Ord a) => a -> a -> a -> Bool}
  \item \mintinline{haskell}{possuiRaizesReais :: Ord a => a -> a -> a -> Bool}

\end{enumerate}

\questionTitle{5}{Tipo de resultado de lookup}
\begin{flushright}
  \footnotesize
  \textit{\textbf{6: Estruturas de dados básicas}}\\
  \textit{c06-lookup-associacao-03}\\

\end{flushright}
Ao buscar uma chave em uma lista de associação com \haskellinline{lookup}, por que o resultado é representado com \haskellinline{Maybe}?

\begin{enumerate}
  \item Porque o uso de \haskellinline{lookup} transforma automaticamente qualquer lista em uma enumeração.

  \item Porque a chave pode estar presente ou ausente, e \haskellinline{Maybe} representa explicitamente essas duas possibilidades.

  \item Porque Haskell não permite devolver diretamente valores como \haskellinline{String} ou \haskellinline{Double}.

  \item Porque toda função que opera sobre tuplas em Haskell precisa devolver um \haskellinline{Maybe}.

  \item Porque \haskellinline{lookup} sempre devolve uma lista, mesmo quando encontra apenas um valor.


\end{enumerate}

\questionTitle{6}{Sintaxe de aplicação prefixa}
\begin{flushright}
  \footnotesize
  \textit{\textbf{3: Expressões e definições}}\\
  \textit{c03-aplicacao-notacoes-01}\\

\end{flushright}
Considere a expressão \haskellinline|mod 25 7|. Qual interpretação está correta de acordo com a sintaxe de aplicação de funções em Haskell?

\begin{enumerate}
  \item \haskellinline|mod| recebe um único argumento, a tupla \haskellinline|(25,7)|.

  \item A função \haskellinline|mod| é aplicada a dois argumentos, \haskellinline|25| e \haskellinline|7|, separados por espaços.

  \item A expressão só seria válida se fosse escrita como \haskellinline|mod(25, 7)|.

  \item A expressão aplica \haskellinline|25| como função aos argumentos \haskellinline|mod| e \haskellinline|7|.

  \item A expressão é inválida porque funções em Haskell sempre exigem parênteses em volta dos argumentos.


\end{enumerate}

\questionTitle{7}{If produzindo função}
\begin{flushright}
  \footnotesize
  \textit{\textbf{5: Expressão condicional}}\\
  \textit{c05-if-avaliacao-03}\\

\end{flushright}
Qual é o resultado da expressão abaixo?
\begin{haskellcode}
(if 5 > 3 then (*) else (+)) 4 2
\end{haskellcode}

\begin{enumerate}
  \item \haskellinline|6|

  \item \haskellinline|8|

  \item Ocorre um erro de sintaxe, pois \haskellinline|if| não pode aparecer entre parênteses.

  \item \haskellinline|14|

  \item Ocorre um erro de tipo, pois os ramos do \haskellinline|if| são funções.


\end{enumerate}

\questionTitle{8}{Raiz quadrada inteira}
\begin{flushright}
  \footnotesize
  \textit{\textbf{9: Funções recursivas}}\\
  \textit{c09-projeto-raiz-inteira-01}\\

\end{flushright}
A raiz quadrada inteira de \haskellinline|n| é o maior \haskellinline|r| tal que \haskellinline|r^2 <= n|. Qual estratégia recursiva está alinhada com essa definição?

\begin{enumerate}
  \item Começar em \haskellinline|0| e diminuir o candidato até que ele seja negativo.

  \item Usar \haskellinline|sqrt| e depois converter o resultado para inteiro, pois o exercício exige ponto flutuante.

  \item Começar com um candidato \haskellinline|i| e diminuí-lo enquanto \haskellinline|i^2 > n|; quando essa condição falhar, \haskellinline|i| é a resposta.

  \item Somar \haskellinline|n| consigo mesmo até ultrapassar \haskellinline|n^2|.

  \item Multiplicar \haskellinline|n| por \haskellinline|2| até que o resultado seja quadrado perfeito.


\end{enumerate}

\questionTitle{9}{Rastreamento da multiplicação recursiva}
\begin{flushright}
  \footnotesize
  \textit{\textbf{9: Funções recursivas}}\\
  \textit{c09-rastreamento-mul-01}\\

\end{flushright}
Considere:
\begin{haskellcode}
mul :: Int -> Int -> Int
mul m n
  | n == 0    = 0
  | n > 0     = m + mul m (n - 1)
  | otherwise = negate (mul m (negate n))
\end{haskellcode}
Qual é o valor de \haskellinline|mul 5 6|?

\begin{enumerate}
  \item \haskellinline|35|

  \item \haskellinline|25|

  \item \haskellinline|30|

  \item \haskellinline|-30|

  \item \haskellinline|11|


\end{enumerate}

\questionTitle{10}{Anotação de tipo em read}
\begin{flushright}
  \footnotesize
  \textit{\textbf{8: Programas interativos}}\\
  \textit{c08-conversao-validacao-03}\\

\end{flushright}
Por que é comum escrever uma anotação como \haskellinline|read linha :: Double| ao converter uma entrada textual?

\begin{enumerate}
  \item Porque \haskellinline|read| é polimórfica. A anotação informa ao compilador qual tipo deve ser produzido pela conversão.

  \item Porque \haskellinline|read| só funciona se o tipo final for escrito dentro de uma string.

  \item Porque \haskellinline|read| sempre retorna \haskellinline|Double|, mas a anotação melhora a legibilidade.

  \item Porque a anotação força o programa a tratar exceções de conversão automaticamente.

  \item Porque sem anotação \haskellinline|read| lê diretamente da entrada padrão.


\end{enumerate}

\questionTitle{11}{Otherwise fora da última guarda}
\begin{flushright}
  \footnotesize
  \textit{\textbf{5: Expressão condicional}}\\
  \textit{c05-otherwise-03}\\

\end{flushright}
Considere a definição:
\begin{haskellcode}
g x
  | otherwise = 1
  | x > 0     = 2
\end{haskellcode}
O que acontece ao avaliar \haskellinline|g 5|?

\begin{enumerate}
  \item Ocorre um erro em tempo de execução por ambiguidade entre as guardas.

  \item O resultado é \haskellinline|2|, porque a guarda \haskellinline|x > 0| seria mais específica.

  \item Ocorre um erro de sintaxe, pois \haskellinline|otherwise| só pode aparecer na última linha.

  \item O compilador reordena as guardas e produz \haskellinline|2|.

  \item O resultado é \haskellinline|1|, porque \haskellinline|otherwise| equivale a \haskellinline|True|.


\end{enumerate}

\questionTitle{12}{Equivalência com guardas}
\begin{flushright}
  \footnotesize
  \textit{\textbf{5: Expressão condicional}}\\
  \textit{c05-if-aninhamento-03}\\

\end{flushright}
Qual definição com guardas é funcionalmente equivalente à função abaixo?
\begin{haskellcode}
sinal n = if n < 0 then -1 else if n > 0 then 1 else 0
\end{haskellcode}

\begin{enumerate}
  \item \begin{haskellcode}
sinal n = if n == 0 then -1 else 1
\end{haskellcode}

  \item \begin{haskellcode}
sinal n
  | n <= 0    = -1
  | otherwise = 1
\end{haskellcode}

  \item \begin{haskellcode}
sinal n
  | n < 0     = -1
  | n > 0     = 1
  | otherwise = 0
\end{haskellcode}

  \item \begin{haskellcode}
sinal n
  | n > 0     = -1
  | n < 0     = 1
  | otherwise = 0
\end{haskellcode}

  \item \begin{haskellcode}
sinal n
  | n < 0     = -1
  | otherwise = 1
\end{haskellcode}


\end{enumerate}

\questionTitle{13}{Rastreamento do acumulador na leitura até zero}
\begin{flushright}
  \footnotesize
  \textit{\textbf{10: Ações de E/S recursivas}}\\
  \textit{c10-acumuladores-cauda-03}\\

\end{flushright}
Na função abaixo, o usuário digita a sequência \haskellinline|10, -2, 7, 0|, um valor por linha.

\begin{haskellcode}
lerESomar :: Integer -> IO Integer
lerESomar total = do
  n <- readLn
  if n == 0
    then return total
    else lerESomar (total + n)
\end{haskellcode}

Qual é o valor de \haskellinline|total| no momento em que o caso base é alcançado?

\begin{enumerate}
  \item 7

  \item 15

  \item 14

  \item 16

  \item 0


\end{enumerate}

\questionTitle{14}{Uso seguro de fromJust}
\begin{flushright}
  \footnotesize
  \textit{\textbf{6: Estruturas de dados básicas}}\\
  \textit{c06-maybe-operacoes-02}\\

\end{flushright}
Em qual situação o uso de \haskellinline{fromJust} é mais justificável neste capítulo?

\begin{enumerate}
  \item Sempre que a expressão aparecer dentro de uma lista.

  \item Quando uma verificação anterior, como \haskellinline{isJust}, já garantiu que o valor não é \haskellinline{Nothing}.

  \item Apenas quando o conteúdo for uma string, mas não quando for um número.

  \item Sempre que a expressão tiver o tipo \haskellinline{Maybe Int}.

  \item Sempre que quisermos evitar o uso de \haskellinline{fromMaybe}.


\end{enumerate}

\questionTitle{15}{Bounded e tipos com limites}
\begin{flushright}
  \footnotesize
  \textit{\textbf{7: Polimorfismo}}\\
  \textit{cap07-hierarquia-bounded-integer}\\

\end{flushright}
Qual alternativa explica corretamente por que \haskellinline{Int} é instância de \haskellinline{Bounded}, mas \haskellinline{Integer} normalmente não é?
\begin{enumerate}
  \item \haskellinline{Int} tem limites mínimo e máximo dependentes da implementação; \haskellinline{Integer} representa inteiros de precisão arbitrária e não possui limite superior ou inferior fixo.
  \item \haskellinline{Integer} não é numérico, por isso não pode ter limites.
  \item \haskellinline{Int} é instância de \haskellinline{Bounded} porque possui valores fracionários.
  \item \haskellinline{Integer} não é instância de \haskellinline{Bounded} porque não é ordenável.
  \item \haskellinline{Bounded} só se aplica a tipos booleanos.

\end{enumerate}

\questionTitle{16}{Tipo de uma média com fromIntegral}
\begin{flushright}
  \footnotesize
  \textit{\textbf{7: Polimorfismo}}\\
  \textit{cap07-expressao-media-corrigida}\\

\end{flushright}
Qual é o tipo mais geral da expressão abaixo?

\begin{minted}{haskell}
fromIntegral (length [True, False]) / 2
\end{minted}

\begin{enumerate}
  \item \mintinline{haskell}{Integer}
  \item \mintinline{haskell}{Bool}
  \item \mintinline{haskell}{Int}
  \item \mintinline{haskell}{[Bool] -> Int}
  \item \mintinline{haskell}{Fractional a => a}

\end{enumerate}

\questionTitle{17}{Erro de tipo ao retornar uma ação}
\begin{flushright}
  \footnotesize
  \textit{\textbf{10: Ações de E/S recursivas}}\\
  \textit{c10-return-tipos-03}\\

\end{flushright}
Considere uma tentativa de escrever uma contagem regressiva recursiva:

\begin{haskellcode}
contagem :: Int -> IO ()
contagem n =
  if n <= 0
    then return ()
    else do
      print n
      return (contagem (n - 1))
\end{haskellcode}

Qual é o problema principal dessa definição?

\begin{enumerate}
  \item A condição \haskellinline|n <= 0| deveria ser substituída por \haskellinline|n == 0|, pois guardas não aceitam desigualdades.

  \item O problema é que ações de tipo \haskellinline|IO ()| nunca podem ser definidas recursivamente.

  \item A chamada \haskellinline|print n| precisa ficar depois da chamada recursiva para que o programa compile.

  \item A função \haskellinline|return| imprime a ação \haskellinline|contagem (n - 1)| em vez de chamá-la.

  \item No ramo recursivo, \haskellinline|return (contagem (n - 1))| cria uma ação que produz outra ação, levando a um tipo como \haskellinline|IO (IO ())| em vez de executar a chamada recursiva.


\end{enumerate}

\questionTitle{18}{Estratégias de Redução e Confluência}
\begin{flushright}
  \footnotesize
  \textit{\textbf{1: Paradigmas de programação}}\\
  \textit{c01-mecanismo-reducao-03}\\

\end{flushright}
Ao avaliar uma expressão matemática em Haskell, o sistema pode teoricamente adotar diferentes caminhos de avaliação, como a redução \emph{de dentro para fora} (aplicativa) ou \emph{de fora para dentro} (normal/preguiçosa). 

Graças à pureza e à ausência de efeitos colaterais na linguagem, qual garantia teórica fundamental o programador possui independentemente do caminho escolhido pelo compilador?

\begin{enumerate}
  \item A garantia de que a quantidade de memória RAM alocada dinamicamente durante o processo de avaliação será anulada pelo \emph{Garbage Collector}.

  \item A garantia de que a execução da expressão consumirá exatamente a mesma quantidade de milissegundos na CPU do usuário.

  \item A garantia de que, caso ambas as estratégias cheguem a um resultado final sem entrar em \emph{loop} infinito, este resultado será sempre o mesmo (Teorema de Church-Rosser).

  \item A garantia de que as funções envolvidas na expressão serão automaticamente vetorizadas e enviadas para execução na GPU.

  \item A garantia de que os tipos numéricos fracionários estarão imunes a qualquer tipo de erro de precisão de ponto flutuante do padrão IEEE 754.


\end{enumerate}

\questionTitle{19}{Definição local em programa interativo}
\begin{flushright}
  \footnotesize
  \textit{\textbf{8: Programas interativos}}\\
  \textit{c08-programas-completos-03}\\

\end{flushright}
Analise:

\begin{minted}{haskell}
main :: IO ()
main = do
  n <- readLn :: IO Int
  let dobro = 2 * n
  print dobro
\end{minted}

Qual afirmação descreve corretamente a linha com \haskellinline|let|?

\begin{enumerate}
  \item Ela define localmente o valor puro \haskellinline|dobro|, usando o valor \haskellinline|n| já obtido pela ação \haskellinline|readLn|.

  \item Ela transforma \haskellinline|dobro| em uma ação do tipo \haskellinline|IO Int|.

  \item Ela só seria válida se fosse escrita como \haskellinline|dobro <- 2 * n|.

  \item Ela executa uma nova ação de E/S para ler o valor de \haskellinline|dobro|.

  \item Ela imprime automaticamente o dobro antes da linha \haskellinline|print dobro|.


\end{enumerate}

\questionTitle{20}{Ambiguidade com read dentro de length}
\begin{flushright}
  \footnotesize
  \textit{\textbf{7: Polimorfismo}}\\
  \textit{cap07-defaulting-length-read}\\

\end{flushright}
Considere:

\begin{minted}{haskell}
r = length [read "10", read "20"]
\end{minted}

Por que essa definição é problemática em compilação normal?

\begin{enumerate}
  \item O resultado de \haskellinline{length} é \haskellinline{Int}, mas o tipo dos elementos lidos permanece restrito apenas por \haskellinline{Read a}, sem contexto suficiente para escolher \haskellinline{a}.
  \item \haskellinline{length} exige que os elementos da lista sejam instâncias de \haskellinline{Num}.
  \item \haskellinline{read} sempre retorna \haskellinline{String}, e \haskellinline{length} não aceita listas de strings.
  \item O problema é que listas em Haskell não podem conter valores produzidos por funções.
  \item A definição é sempre resolvida como \haskellinline{[Int]} por causa dos caracteres numéricos nas strings.

\end{enumerate}

\questionTitle{21}{Tipo unitário}
\begin{flushright}
  \footnotesize
  \textit{\textbf{6: Estruturas de dados básicas}}\\
  \textit{c06-tuplas-unitario-01}\\

\end{flushright}
Qual afirmação descreve corretamente o tipo \haskellinline{()} em Haskell?

\begin{enumerate}
  \item É apenas outra notação para a lista vazia \haskellinline{[]}.

  \item É o tipo unitário: possui exatamente um único valor, também escrito como \haskellinline{()}.

  \item É a forma como Haskell representa uma tupla de um elemento.

  \item É um tipo reservado exclusivamente para expressões numéricas.

  \item É um sinônimo de tipo para \haskellinline{Bool}.


\end{enumerate}

\questionTitle{22}{Resultados de funções RealFrac}
\begin{flushright}
  \footnotesize
  \textit{\textbf{7: Polimorfismo}}\\
  \textit{cap07-conversao-realfrac-param}\\

\end{flushright}
Qual é o resultado da avaliação da expressão abaixo?

\begin{minted}{haskell}
floor (-2.8) :: Int
\end{minted}

\begin{enumerate}
  \item A expressão causa erro de tipo, pois \haskellinline{RealFrac} não permite retornar \haskellinline{Int}.
  \item \haskellinline|2|
  \item \haskellinline|-3|
  \item \haskellinline|-2|
  \item \haskellinline|3|

\end{enumerate}

\questionTitle{23}{Garantias da função pura}
\begin{flushright}
  \footnotesize
  \textit{\textbf{1: Paradigmas de programação}}\\
  \textit{c01-transparencia-referencial-03}\\

\end{flushright}
Uma \textbf{função pura} é a base fundamental da construção de software no paradigma declarativo funcional. Assinale a alternativa que descreve corretamente as duas garantias absolutas oferecidas por uma função pura.

\begin{enumerate}
  \item 1) Ela é garantida de terminar (não entrar em loop infinito). 2) Ela executa obrigatoriamente em tempo constante $\mathcal{O}(1)$.

  \item 1) Ela avalia todos os seus argumentos de forma estrita (\emph{strict evaluation}) antes de executar. 2) Ela não depende de bibliotecas de terceiros.

  \item 1) Ela possui uma assinatura de tipo polimórfica explícita. 2) Ela substitui automaticamente operações de ponto flutuante por inteiros para evitar imprecisões matemáticas.

  \item 1) Ela não utiliza declarações de variáveis locais em seu corpo. 2) Ela só pode receber matrizes e listas imutáveis como parâmetros.

  \item 1) Ela sempre avalia para o mesmo resultado quando recebe os mesmos argumentos. 2) Ela não causa nenhum efeito colateral observável fora de seu escopo.


\end{enumerate}

\questionTitle{24}{Ações de saída padrão}
\begin{flushright}
  \footnotesize
  \textit{\textbf{8: Programas interativos}}\\
  \textit{c08-entrada-saida-03}\\

\end{flushright}
Considere a função de saída \haskellinline|print|. Qual alternativa descreve corretamente seu tipo e seu comportamento?

\begin{enumerate}
  \item \haskellinline|String -> String| — transforma texto sem construir nenhuma ação de E/S.

  \item \haskellinline|String -> IO ()| — só imprime strings já formatadas.

  \item \haskellinline|Read a => IO a| — lê e converte valores da entrada padrão.

  \item \haskellinline|Show a => a -> IO ()| — constrói uma ação que imprime a representação textual de um valor.

  \item \haskellinline|IO String| — lê da entrada padrão o texto que deveria ser impresso.


\end{enumerate}

\questionTitle{25}{Avaliação parametrizada com precedência e associatividade}
\begin{flushright}
  \footnotesize
  \textit{\textbf{3: Expressões e definições}}\\
  \textit{c03-precedencia-associatividade-01}\\

\end{flushright}
Qual é o resultado da seguinte expressão em Haskell?

\begin{haskellcode}
length [1,2,3] + 4 * 2
\end{haskellcode}

\begin{enumerate}
  \item A expressão é sintaticamente inválida.

  \item \haskellinline|14|.

  \item \haskellinline|11|.

  \item \haskellinline|8|.

  \item \haskellinline|9|.


\end{enumerate}

\questionTitle{26}{Necessidade de reverse ao acumular com cons}
\begin{flushright}
  \footnotesize
  \textit{\textbf{10: Ações de E/S recursivas}}\\
  \textit{c10-listas-ordem-01}\\

\end{flushright}
Considere a função:

\begin{haskellcode}
lerLista :: [Integer] -> IO [Integer]
lerLista xs = do
  x <- readLn
  if x == 0
    then return (reverse xs)
    else lerLista (x:xs)
\end{haskellcode}

Por que \haskellinline|reverse| é usado no caso base?

\begin{enumerate}
  \item Porque \haskellinline|reverse| força a ação de E/S a ser executada antes do caso base.

  \item Porque \haskellinline|sum| exige que a lista esteja invertida para funcionar corretamente.

  \item Porque \haskellinline|readLn| sempre lê os números na ordem inversa em que foram digitados.

  \item Porque cada novo valor é inserido no início do acumulador com \haskellinline|x:xs|, invertendo a ordem de entrada.

  \item Porque \haskellinline|return| só aceita listas ordenadas em ordem crescente.


\end{enumerate}

\questionTitle{27}{O comando :type e sua utilidade}
\begin{flushright}
  \footnotesize
  \textit{\textbf{2: Ambiente de desenvolvimento}}\\
  \textit{c02-ghci-inspecao-01}\\

\end{flushright}
Durante uma sessão interativa no GHCi, o programador utiliza frequentemente o comando \texttt{:type} (ou \texttt{:t}). O que este comando executa quando aplicado a uma expressão?

\begin{enumerate}
  \item Ele exibe o código-fonte original (em linguagem Haskell) que originou a função solicitada.

  \item Ele força a avaliação completa da expressão matemática e formata a saída em um texto de múltiplas linhas legível para o usuário.

  \item Ele modifica o tipo da variável em tempo real, permitindo que um inteiro seja tratado como \emph{string} no GHCi temporariamente.

  \item Ele pesquisa no repositório oficial da linguagem quais pacotes exportam tipos de dados que correspondam à expressão digitada.

  \item Ele inspeciona e exibe a assinatura de tipo inferida pelo compilador para a expressão fornecida, sem efetivamente avaliá-la (executá-la).


\end{enumerate}

\questionTitle{28}{Classes de tipo e interfaces de POO}
\begin{flushright}
  \footnotesize
  \textit{\textbf{7: Polimorfismo}}\\
  \textit{cap07-comparacao-typeclass-poo}\\

\end{flushright}
Uma analogia comum aproxima classes de tipo em Haskell de interfaces em linguagens orientadas a objetos. Qual afirmação evita a simplificação excessiva?
\begin{enumerate}
  \item A analogia é útil porque ambas especificam operações disponíveis, mas classes de tipo são resolvidas por instâncias e restrições de tipos, não por subtipagem nominal de objetos mutáveis.
  \item Classes de tipo só existem para simular herança múltipla em Haskell.
  \item Classes de tipo não têm qualquer relação com operações; elas servem apenas para documentação.
  \item Classes de tipo são exatamente classes de POO, com campos, construtores e herança de implementação.
  \item Interfaces de POO e classes de tipo são idênticas porque ambas exigem que todos os métodos sejam chamados dinamicamente.

\end{enumerate}

\questionTitle{29}{Suficiência do tipo Int}
\begin{flushright}
  \footnotesize
  \textit{\textbf{4: Tipos de Dados}}\\
  \textit{c04-tipos-int-integer-02}\\

\end{flushright}
Para qual dos seguintes valores o tipo \haskellinline{Int} é geralmente suficiente, não necessitando de \haskellinline{Integer}?

\begin{enumerate}
  \item A distância em milímetros entre duas galáxias.

  \item A idade de uma pessoa.

  \item O número de grãos de areia na Terra.

  \item O resultado de \haskellinline{2 ^ 200}.

  \item A população mundial.


\end{enumerate}

\questionTitle{30}{Erro por desalinhamento na regra de layout}
\begin{flushright}
  \footnotesize
  \textit{\textbf{3: Expressões e definições}}\\
  \textit{c03-layout-avaliacao-01}\\

\end{flushright}
Considere o trecho abaixo, no qual a segunda definição local não está alinhada com a primeira.

\begin{haskellcode}
f x = let a = x + 1
        b = x + 2
      in a + b
\end{haskellcode}

Qual é a causa do problema de sintaxe?

\begin{enumerate}
  \item O problema é que \haskellinline|in| deve aparecer antes de todas as definições do bloco.

  \item Pela regra de layout, definições pertencentes ao mesmo bloco devem começar na mesma coluna; \haskellinline|b| deveria estar alinhado com \haskellinline|a|.

  \item O problema é o uso de \haskellinline|let|; definições locais só podem ser feitas com \haskellinline|where|.

  \item O problema é que Haskell ignora espaços em branco, então o compilador não consegue separar tokens.

  \item O problema é o uso de \haskellinline|+| dentro de definições locais.


\end{enumerate}

\questionTitle{31}{Paralelismo Seguro}
\begin{flushright}
  \footnotesize
  \textit{\textbf{1: Paradigmas de programação}}\\
  \textit{c01-engenharia-software-03}\\

\end{flushright}
No contexto de sistemas de alta disponibilidade executados em processadores \emph{multicore}, a programação funcional tem ganhado vasto espaço na indústria contemporânea. Qual característica fundacional do paradigma facilita imensamente a construção de programas paralelos livres de falhas?

\begin{enumerate}
  \item A tipagem dinâmica fraca da linguagem assegura que os processos em segundo plano reescrevam as estruturas de memória de forma atômica e sincronizada por meio do \emph{garbage collector}.

  \item O paradigma impossibilita deliberadamente a criação de processos paralelos (sendo restrito a \emph{single-thread}), o que mitiga os erros de concorrência anulando a complexidade do escalonamento.

  \item Os compiladores de linguagens funcionais traduzem nativamente o código para chamadas assíncronas de \emph{hardware} que contornam totalmente os barramentos lentos da memória principal.

  \item A avaliação estrita das linguagens funcionais impõe um controle rígido ao escalonador do sistema operacional, forçando a distribuição milimétrica da carga da aplicação entre os núcleos físicos.

  \item A imutabilidade dos dados combinada às funções puras garantem que múltiplas \emph{threads} não causarão \emph{race conditions} (condições de corrida) tentando modificar o mesmo espaço de memória simultaneamente.


\end{enumerate}

\questionTitle{32}{Compatibilidade dos ramos}
\begin{flushright}
  \footnotesize
  \textit{\textbf{5: Expressão condicional}}\\
  \textit{c05-if-regras-03}\\

\end{flushright}
Qual expressão abaixo é rejeitada porque os ramos \haskellinline|then| e \haskellinline|else| não produzem valores do mesmo tipo?

\begin{enumerate}
  \item \haskellinline|if odd 5 then 1 else 0|

  \item \haskellinline|if odd 5 then 'i' else 'p'|

  \item \haskellinline|if odd 5 then [1,2] else [3,4]|

  \item \haskellinline|if odd 5 then "ímpar" else "par"|

  \item \haskellinline|if odd 5 then "ímpar" else 0|


\end{enumerate}

\questionTitle{33}{Guarda sem cobertura completa}
\begin{flushright}
  \footnotesize
  \textit{\textbf{9: Funções recursivas}}\\
  \textit{c09-fundamentos-cobertura-01}\\

\end{flushright}
Considere a definição:
\begin{haskellcode}
h :: Integer -> Integer
h n
  | n == 0 = 1
  | n > 0  = n * h (n - 1)
\end{haskellcode}
Qual alternativa descreve corretamente o comportamento de \haskellinline|h (-2)|?

\begin{enumerate}
  \item Nenhuma guarda é satisfeita para \haskellinline|(-2)|; portanto, a avaliação falha por ausência de caso aplicável.

  \item A avaliação entra em recursão infinita, pois a guarda \haskellinline|n > 0| também vale para negativos.

  \item O resultado é \haskellinline|1|, pois a primeira guarda é usada para qualquer valor não positivo.

  \item O resultado é \haskellinline|-2|, pois a função devolve o argumento quando ele é negativo.

  \item Ocorre erro de tipo, pois guardas booleanas não podem ser usadas com \haskellinline|Integer|.


\end{enumerate}

\questionTitle{34}{Inferência de tipo de função composta}
\begin{flushright}
  \footnotesize
  \textit{\textbf{4: Tipos de Dados}}\\
  \textit{c04-inferencia-funcao-01}\\

\end{flushright}
Considere a função definida por \haskellinline{f x = not (x > 0)}. Qual é o tipo inferido para \haskellinline{f}?

\begin{enumerate}
  \item \haskellinline{Int -> Bool}

  \item \haskellinline{(Ord a, Num a) => a -> Bool}

  \item \haskellinline{a -> Bool}

  \item \haskellinline{Num a => a -> Bool}

  \item \haskellinline{Double -> Bool}


\end{enumerate}

\questionTitle{35}{Assinatura em expressões}
\begin{flushright}
  \footnotesize
  \textit{\textbf{4: Tipos de Dados}}\\
  \textit{c04-assinatura-expressao-01}\\

\end{flushright}
Por que a expressão \haskellinline{read "42"} sozinha causaria um erro de ambiguidade de tipo, enquanto \haskellinline{read "42" :: Int} é aceita?

\begin{enumerate}
  \item A anotação \haskellinline{:: Int} força a conversão da string para inteiro em tempo de execução.

  \item \haskellinline{read} é sobrecarregada e pode retornar vários tipos; a anotação \haskellinline{:: Int} resolve a ambiguidade.

  \item A anotação é necessária porque \haskellinline{read} é uma função parcial.

  \item \haskellinline{read} só funciona com anotação de tipo; sem ela, ocorre um erro de sintaxe.

  \item A anotação \haskellinline{:: Int} instrui o compilador a ignorar a inferência de tipos.


\end{enumerate}

\questionTitle{36}{Análise de IMC com where}
\begin{flushright}
  \footnotesize
  \textit{\textbf{5: Expressão condicional}}\\
  \textit{c05-where-avaliacao-03}\\

\end{flushright}
Considere a função:
\begin{haskellcode}
analisaIMC peso altura
  | imc < 18.5 = "Abaixo"
  | imc < 25.0 = "Normal"
  | otherwise  = "Acima"
  where
    imc = peso / (altura * altura)
\end{haskellcode}
Qual é o resultado de \haskellinline|analisaIMC 80 2.0|?

\begin{enumerate}
  \item \haskellinline|"80.0"|

  \item \haskellinline|"Abaixo"|

  \item Ocorre um erro, pois variáveis definidas em \haskellinline|where| não podem aparecer nas guardas.

  \item \haskellinline|"Normal"|

  \item \haskellinline|"Acima"|


\end{enumerate}

\questionTitle{37}{Cons versus concatenação}
\begin{flushright}
  \footnotesize
  \textit{\textbf{6: Estruturas de dados básicas}}\\
  \textit{c06-listas-estrutura-03}\\

\end{flushright}
Qual afirmação descreve corretamente a diferença entre \haskellinline{(:)} e \haskellinline{(++)}?

\begin{enumerate}
  \item \haskellinline{(++)} adiciona um elemento à lista, enquanto \haskellinline{(:)} exige duas listas como argumentos.

  \item \haskellinline{(:)} adiciona um elemento ao início de uma lista, enquanto \haskellinline{(++)} concatena duas listas.

  \item \haskellinline{(:)} adiciona um elemento ao final da lista, enquanto \haskellinline{(++)} adiciona ao início.

  \item \haskellinline{(:)} e \haskellinline{(++)} sempre produzem exatamente o mesmo efeito, mudando apenas a sintaxe.

  \item \haskellinline{(:)} concatena listas mais rapidamente do que \haskellinline{(++)}, mas exige listas do mesmo tamanho.


\end{enumerate}

\questionTitle{38}{Finalidade da cláusula where}
\begin{flushright}
  \footnotesize
  \textit{\textbf{3: Expressões e definições}}\\
  \textit{c03-definicoes-where-01}\\

\end{flushright}
Qual alternativa descreve corretamente o papel da cláusula \haskellinline|where| em uma definição Haskell?

\begin{enumerate}
  \item Ela introduz uma expressão que pode ser escrita livremente em qualquer posição onde um literal seria permitido.

  \item Ela serve exclusivamente para importar nomes de outros módulos.

  \item Ela anexa definições locais a uma equação, permitindo nomear subexpressões auxiliares no escopo dessa definição.

  \item Ela torna todas as definições auxiliares visíveis globalmente no módulo.

  \item Ela substitui a necessidade do sinal de igualdade em definições de funções.


\end{enumerate}

\questionTitle{39}{Média de sequência vazia}
\begin{flushright}
  \footnotesize
  \textit{\textbf{10: Ações de E/S recursivas}}\\
  \textit{c10-bordas-robustez-03}\\

\end{flushright}
Considere o padrão usado no exercício da média:

\begin{haskellcode}
main :: IO ()
main = do
  lista <- lerLista
  print (sum lista / fromIntegral (length lista))

lerLista :: IO [Double]
lerLista = do
  x <- readLn
  if x < 0
    then return []
    else do
      xs <- lerLista
      return (x:xs)
\end{haskellcode}

Qual caso de borda precisa ser tratado explicitamente?

\begin{enumerate}
  \item A lista com um único valor, pois \haskellinline|sum| não funciona com listas unitárias.

  \item A presença de zero na lista, pois zero é sempre interpretado como sentinela nessa função.

  \item Qualquer lista com valores positivos, pois \haskellinline|fromIntegral| só aceita números negativos.

  \item A lista em ordem crescente, pois \haskellinline|sum| exige que a entrada esteja ordenada.

  \item A sequência vazia, pois se o primeiro valor digitado for negativo, \haskellinline|length lista| será zero e a média não estará bem definida.


\end{enumerate}

\questionTitle{40}{Paradigmas - Característica Fundamental (Legado Refatorado)}
\begin{flushright}
  \footnotesize
  \textit{\textbf{1: Paradigmas de programação}}\\
  \textit{c01-conceitos-paradigmas-01}\\

\end{flushright}
Qual das seguintes características é a mais fundamental para definir uma linguagem como \textbf{puramente funcional}?

\begin{enumerate}
  \item O tratamento de funções como valores de primeira classe e a total ausência de efeitos colaterais.

  \item O uso de ponteiros abstratos para a manipulação direta e manual da memória pelo programador.

  \item A capacidade de usar laços imperativos nativos como \cinline|for| e \cinline|while|.

  \item A ausência completa de um sistema de tipos estático.

  \item A obrigatoriedade de usar o paradigma de orientação a objetos para gerenciar o estado da aplicação.


\end{enumerate}

\questionTitle{41}{Declaração de sinônimo de tipo}
\begin{flushright}
  \footnotesize
  \textit{\textbf{4: Tipos de Dados}}\\
  \textit{c04-sinonimos-criacao-01}\\

\end{flushright}
Qual palavra-chave é usada para criar um sinônimo de tipo em Haskell?

\begin{enumerate}
  \item \haskellinline{data}

  \item \haskellinline{instance}

  \item \haskellinline{type}

  \item \haskellinline{class}

  \item \haskellinline{newtype}


\end{enumerate}

\questionTitle{42}{String vs IO String}
\begin{flushright}
  \footnotesize
  \textit{\textbf{8: Programas interativos}}\\
  \textit{c08-tipos-acoes-03}\\

\end{flushright}
Considere o trecho:

\begin{minted}{haskell}
nome = getLine
\end{minted}

Sem nenhuma outra informação de contexto, qual é o tipo de \haskellinline|nome|?

\begin{enumerate}
  \item \haskellinline|IO String|, pois \haskellinline|getLine| é uma ação de E/S que, quando executada, produz uma \haskellinline|String|.

  \item \haskellinline|Char|, pois a entrada padrão lê apenas um caractere por vez.

  \item \haskellinline|String -> IO String|, pois \haskellinline|getLine| precisa receber o prompt como argumento.

  \item \haskellinline|String|, pois \haskellinline|getLine| já executa a leitura no momento da definição.

  \item \haskellinline|IO ()|, pois toda ação de E/S sempre produz o valor unitário.


\end{enumerate}

\questionTitle{43}{Modelagem composta de pedidos}
\begin{flushright}
  \footnotesize
  \textit{\textbf{6: Estruturas de dados básicas}}\\
  \textit{c06-modelagem-03}\\

\end{flushright}
Qual tipo representa melhor uma coleção de pedidos em que cada pedido possui o nome do cliente, a lista de itens comprados e um total opcional?

\begin{enumerate}
  \item \haskellinline{Maybe [(String, Int, Double)]}

  \item \haskellinline{[(String, [(String, Int)], Maybe Double)]}

  \item \haskellinline{([String], [Int], [Double])}

  \item \haskellinline{[String, [(String, Int)], Maybe Double]}

  \item \haskellinline{(String, [String], Double)}


\end{enumerate}

\questionTitle{44}{Composição de funções paramétricas em listas}
\begin{flushright}
  \footnotesize
  \textit{\textbf{7: Polimorfismo}}\\
  \textit{cap07-param-head-tail-tipo}\\

\end{flushright}
Qual é a assinatura de tipo mais geral da função abaixo?

\begin{minted}{haskell}
f xs = head (tail xs)
\end{minted}

\begin{enumerate}
  \item \mintinline{haskell}{f :: [a] -> [a]}
  \item \mintinline{haskell}{f :: Ord a => [a] -> a}
  \item \mintinline{haskell}{f :: [a] -> a}
  \item \mintinline{haskell}{f :: [[a]] -> [a]}
  \item \mintinline{haskell}{f :: Num a => [a] -> a}

\end{enumerate}

\questionTitle{45}{O ciclo prático de edição e recarregamento}
\begin{flushright}
  \footnotesize
  \textit{\textbf{2: Ambiente de desenvolvimento}}\\
  \textit{c02-fluxo-trabalho-01}\\

\end{flushright}
Qual das seguintes sequências descreve adequadamente o fluxo de trabalho (\emph{workflow}) mais ágil e recomendado para o desenvolvimento de programas locais usando Haskell e GHCi?

\begin{enumerate}
  \item Modificar o arquivo \texttt{.hs} no editor $\rightarrow$ Fechar o GHCi $\rightarrow$ Invocar o compilador GHC na linha de comando para gerar um novo pacote executável $\rightarrow$ Testar o binário.

  \item Editar o código-fonte em um editor externo e salvá-lo $\rightarrow$ Usar o comando \texttt{:reload} (ou \texttt{:r}) no GHCi já aberto $\rightarrow$ Testar a nova função iterativamente.

  \item Digitar o comando \texttt{:load} no editor VS Code $\rightarrow$ Modificar a função $\rightarrow$ Enviar a compilação final para o GitHub Classroom rodar as ferramentas de análise.

  \item Reescrever todo o código no GHCi $\rightarrow$ Digitar o comando de exportação de binário $\rightarrow$ Rodar o \emph{script} nativo em C.

  \item Editar as funções diretamente dentro do GHCi através da declaração múltipla $\rightarrow$ Encerrar o GHCi com \texttt{:quit} $\rightarrow$ Reabrir o terminal para que as alterações sejam gravadas no disco.


\end{enumerate}

\questionTitle{46}{A variável automática it no GHCi}
\begin{flushright}
  \footnotesize
  \textit{\textbf{3: Expressões e definições}}\\
  \textit{c03-ghci-definicoes-01}\\

\end{flushright}
Considere a seguinte sessão no GHCi.

\begin{haskellcode}
> 2 + 3 * 4
14
> 7 * (it - 4)
70
> it
\end{haskellcode}

Qual valor será exibido na última linha?

\begin{enumerate}
  \item \haskellinline|2 + 3 * 4|.

  \item \haskellinline|14|.

  \item \haskellinline|10|.

  \item \haskellinline|70|.

  \item O GHCi produz erro, pois \haskellinline|it| só pode ser usado uma vez.


\end{enumerate}

\questionTitle{47}{Conflito de classes em checagem estática}
\begin{flushright}
  \footnotesize
  \textit{\textbf{4: Tipos de Dados}}\\
  \textit{c04-checagem-conflito-01}\\

\end{flushright}
Considere o código \haskellinline{calculo x y = (x / 2) `div` y}. Por que este código falha na compilação?

\begin{enumerate}
  \item A notação infixa com crases não pode ser usada com operadores simbólicos.

  \item Falta uma anotação de tipo explícita, impedindo a inferência.

  \item O operador \haskellinline{/} exige um tipo fracionário (\haskellinline{Fractional}), mas \haskellinline{div} exige um tipo integral (\haskellinline{Integral}), causando um conflito de tipos.

  \item O literal \haskellinline{2} é inferido como \haskellinline{Integer}, que não é compatível com \haskellinline{x}.

  \item A função \haskellinline{div} tem maior precedência que \haskellinline{/}.


\end{enumerate}

\questionTitle{48}{Identificando tipos vs. variáveis}
\begin{flushright}
  \footnotesize
  \textit{\textbf{4: Tipos de Dados}}\\
  \textit{c04-conceito-nomes-02}\\

\end{flushright}
Considere os identificadores: \haskellinline{Pessoa} e \haskellinline{pessoa}. Qual é a interpretação mais provável em Haskell?

\begin{enumerate}
  \item Ambos são tipos, sendo \haskellinline{pessoa} um sinônimo.

  \item \haskellinline{Pessoa} é um tipo (ou construtor de dados) e \haskellinline{pessoa} é uma variável ou função.

  \item \haskellinline{Pessoa} é uma variável global e \haskellinline{pessoa} é uma variável local.

  \item Ambos são variáveis, mas com escopos diferentes.

  \item Não há diferença; Haskell não distingue maiúsculas de minúsculas.


\end{enumerate}

\questionTitle{49}{Declaração de módulo em arquivo fonte}
\begin{flushright}
  \footnotesize
  \textit{\textbf{3: Expressões e definições}}\\
  \textit{c03-arquivos-identificadores-01}\\

\end{flushright}
Qual alternativa apresenta uma declaração de módulo bem formada para o início de um arquivo Haskell?

\begin{enumerate}
  \item \haskellinline|module Geometria =|.

  \item \haskellinline|import module Geometria where|.

  \item \haskellinline|module geometria where|.

  \item \haskellinline|module Geometria where|.

  \item \haskellinline|where module Geometria|.


\end{enumerate}

\questionTitle{50}{A inovação do Sistema de Tipos da ML}
\begin{flushright}
  \footnotesize
  \textit{\textbf{1: Paradigmas de programação}}\\
  \textit{c01-historia-evolucao-03}\\

\end{flushright}
Na década de 1970, o pesquisador Robin Milner e sua equipe desenvolveram a linguagem ML (\emph{MetaLanguage}), projetada inicialmente para apoiar a demonstração automática de teoremas. A linguagem ML legou um avanço formidável, hoje presente fortemente em Haskell e linguagens como Rust e Swift. Qual foi essa inovação técnica?

\begin{enumerate}
  \item A introdução e formalização das \textbf{Mônadas} na biblioteca padrão como a única abstração autorizada pelo compilador para isolar e manipular efeitos de entrada e saída.

  \item A arquitetura inovadora do compilador \textbf{JIT} (\emph{Just-in-Time Compiler}), que otimizou o código funcional para superar a velocidade dos programas escritos em linguagem C de forma absoluta.

  \item O mecanismo de \textbf{Avaliação Preguiçosa} estrita, que impedia o cálculo precipitado de matrizes infinitas e protegia os \emph{mainframes} contra falhas de estouro de memória.

  \item Um sistema de \textbf{inferência de tipos} elegante e poderoso, capaz de garantir a segurança estática sem exigir que o programador preenchesse o código com anotações manuais redundantes de tipos.

  \item A invenção pioneira da \textbf{Orientação a Objetos}, que permitiu à indústria mesclar funções puras com o encapsulamento hermético de estado dentro de classes herdáveis.


\end{enumerate}

\questionTitle{51}{Restrição Read e necessidade de contexto}
\begin{flushright}
  \footnotesize
  \textit{\textbf{7: Polimorfismo}}\\
  \textit{cap07-adhoc-read-contexto}\\

\end{flushright}
A função \haskellinline{read} tem assinatura:

\begin{minted}{haskell}
read :: Read a => String -> a
\end{minted}

Qual uso deixa claro para o compilador que o resultado deve ser um \haskellinline{Int}?

\begin{enumerate}
  \item \mintinline{haskell}{read "10" :: Int}
  \item \mintinline{haskell}{read "10" :: String -> Int}
  \item \mintinline{haskell}{read :: "10" -> Int}
  \item \mintinline{haskell}{show (read "10")}
  \item \mintinline{haskell}{length (read "10")}

\end{enumerate}

\questionTitle{52}{Classificação por faixas}
\begin{flushright}
  \footnotesize
  \textit{\textbf{5: Expressão condicional}}\\
  \textit{c05-guardas-avaliacao-03}\\

\end{flushright}
Analise a definição:
\begin{haskellcode}
conceito n
  | n >= 9    = "A"
  | n >= 7    = "B"
  | n >= 5    = "C"
  | otherwise = "D"
\end{haskellcode}
Qual é o resultado de \haskellinline|conceito 7.5|?

\begin{enumerate}
  \item \haskellinline|"C"|

  \item \haskellinline|"D"|

  \item \haskellinline|"A"|

  \item Ocorre um erro de tipos, pois as guardas usam números decimais.

  \item \haskellinline|"B"|


\end{enumerate}

\questionTitle{53}{Acumuladores e avaliação preguiçosa}
\begin{flushright}
  \footnotesize
  \textit{\textbf{9: Funções recursivas}}\\
  \textit{c09-eficiencia-thunks-01}\\

\end{flushright}
Em Haskell, por que uma função com acumulador ainda pode consumir mais memória do que o esperado?

\begin{enumerate}
  \item Porque a avaliação preguiçosa pode adiar o cálculo do acumulador, formando expressões suspensas, ou \emph{thunks}.

  \item Porque Haskell não permite que acumuladores sejam avaliados.

  \item Porque funções com acumulador não podem ter casos base.

  \item Porque o tipo \haskellinline|Integer| impede qualquer otimização do compilador.

  \item Porque acumuladores são sempre armazenados como listas infinitas.


\end{enumerate}

\questionTitle{54}{Construção e execução de ações}
\begin{flushright}
  \footnotesize
  \textit{\textbf{8: Programas interativos}}\\
  \textit{c08-modelo-io-03}\\

\end{flushright}
Considere a definição:

\begin{minted}{haskell}
saudacao :: IO ()
saudacao = putStrLn "Olá!"
\end{minted}

O que essa definição faz no momento em que o módulo é carregado?

\begin{enumerate}
  \item Ela executa a ação uma vez e guarda o resultado textual \haskellinline|"Olá!"| para reutilização futura.

  \item Ela define uma função pura do tipo \haskellinline|String -> String| que transforma a mensagem antes de imprimi-la.

  \item Ela cria uma exceção, pois \haskellinline|putStrLn| só pode aparecer diretamente dentro de \haskellinline|main|.

  \item Ela imprime imediatamente \haskellinline|"Olá!"|, pois toda definição de tipo \haskellinline|IO ()| é executada ao carregar o módulo.

  \item Ela associa o nome \haskellinline|saudacao| a uma ação de E/S. A mensagem só será impressa quando essa ação for executada.


\end{enumerate}

\questionTitle{55}{Papel de where em guardas}
\begin{flushright}
  \footnotesize
  \textit{\textbf{5: Expressão condicional}}\\
  \textit{c05-where-escopo-03}\\

\end{flushright}
Qual é a principal vantagem de usar \haskellinline|where| em uma definição com guardas, como no cálculo de \haskellinline|delta| em uma função que determina o número de raízes de uma equação?

\begin{enumerate}
  \item Tornar a função automaticamente polimórfica.

  \item Garantir que todas as guardas sejam exaustivas.

  \item Permitir nomear um cálculo local e reutilizá-lo em guardas e resultados, evitando repetição.

  \item Permitir que a função tenha ramos com tipos diferentes.

  \item Fazer com que as guardas sejam avaliadas em paralelo.


\end{enumerate}

\questionTitle{56}{let não executa getLine}
\begin{flushright}
  \footnotesize
  \textit{\textbf{8: Programas interativos}}\\
  \textit{c08-blocos-do-03}\\

\end{flushright}
Considere:

\begin{minted}{haskell}
saudacao :: IO ()
saudacao = do
  putStrLn "Qual o seu nome?"
  let nome = getLine
  putStrLn "Olá!"
\end{minted}

O que acontece quando \haskellinline|saudacao| é executada?

\begin{enumerate}
  \item O programa não compila, pois \haskellinline|let| é sintaticamente proibido dentro de \haskellinline|do|.

  \item O programa imprime a pergunta, aguarda o nome e imprime \haskellinline|"Olá!"|.

  \item O programa imprime a pergunta e depois imprime \haskellinline|"Olá!"|, mas não aguarda a digitação do nome. A linha com \haskellinline|let| apenas associa \haskellinline|nome| à ação \haskellinline|getLine|.

  \item O programa imprime a ação \haskellinline|getLine| como texto antes de imprimir \haskellinline|"Olá!"|.

  \item O programa entra em laço infinito porque \haskellinline|getLine| fica esperando entrada indefinidamente.


\end{enumerate}

\questionTitle{57}{Saída produzida por mostraLista}
\begin{flushright}
  \footnotesize
  \textit{\textbf{10: Ações de E/S recursivas}}\\
  \textit{c10-rastreamento-03}\\

\end{flushright}
Considere:

\begin{haskellcode}
mostraLista :: Show a => [a] -> IO ()
mostraLista lista
  | null lista = return ()
  | otherwise = do
      print (head lista)
      mostraLista (tail lista)
\end{haskellcode}

Qual é a saída de \haskellinline|mostraLista [2,4,6]|?

\begin{enumerate}
  \item Apenas \haskellinline|2| é impresso, pois a chamada recursiva é ignorada.

  \item Os valores \haskellinline|6|, \haskellinline|4| e \haskellinline|2| são impressos, um por linha.

  \item Nada é impresso, pois a função retorna \haskellinline|()|.

  \item \haskellinline|[2,4,6]| em uma única linha.

  \item Os valores \haskellinline|2|, \haskellinline|4| e \haskellinline|6| são impressos, um por linha, nessa ordem.


\end{enumerate}

\questionTitle{58}{GHCup e gerenciamento de ambiente}
\begin{flushright}
  \footnotesize
  \textit{\textbf{2: Ambiente de desenvolvimento}}\\
  \textit{c02-ecossistema-ferramentas-01}\\

\end{flushright}
No ecossistema moderno da linguagem Haskell, qual é o papel exato da ferramenta \textbf{GHCup}?

\begin{enumerate}
  \item Fornecer a interface gráfica e o editor de texto oficial recomendado pela comunidade para programação Haskell.

  \item Converter automaticamente o código-fonte escrito em arquivos \texttt{.hs} para binários otimizados da linguagem C.

  \item Servir como o interpretador \emph{REPL} principal para testar e avaliar expressões matemáticas interativamente.

  \item Hospedar repositórios de código-fonte na nuvem, similar ao GitHub, mas com foco exclusivo em bibliotecas funcionais.

  \item Atuar como um instalador unificado que gerencia as versões do compilador GHC, da ferramenta de \emph{build} Cabal e do servidor de linguagem HLS.


\end{enumerate}

\questionTitle{59}{Guardas sobrepostas}
\begin{flushright}
  \footnotesize
  \textit{\textbf{5: Expressão condicional}}\\
  \textit{c05-guardas-sintaxe-semantica-03}\\

\end{flushright}
Considere a função:
\begin{haskellcode}
f x
  | x > 0     = 1
  | x > 10    = 2
  | otherwise = 3
\end{haskellcode}
Qual é o valor de \haskellinline|f 20|?

\begin{enumerate}
  \item O compilador rejeita a função, pois as guardas se sobrepõem.

  \item Ocorre um erro, pois duas guardas são verdadeiras.

  \item \haskellinline|3|

  \item \haskellinline|2|

  \item \haskellinline|1|


\end{enumerate}

\questionTitle{60}{Rastreamento do fatorial com acumulador}
\begin{flushright}
  \footnotesize
  \textit{\textbf{9: Funções recursivas}}\\
  \textit{c09-cauda-fat-01}\\

\end{flushright}
Considere:
\begin{haskellcode}
fat' n x
  | n == 0 = x
  | n > 0  = fat' (n - 1) (n * x)
\end{haskellcode}
Qual é a próxima chamada feita por \haskellinline|fat' 5 2|?

\begin{enumerate}
  \item \haskellinline|fat' 4 2|

  \item \haskellinline|fat' 10 4|

  \item \haskellinline|fat' 4 10|

  \item \haskellinline|fat' 0 240|

  \item \haskellinline|fat' 5 10|


\end{enumerate}

\questionTitle{61}{NoBuffering}
\begin{flushright}
  \footnotesize
  \textit{\textbf{8: Programas interativos}}\\
  \textit{c08-buffering-03}\\

\end{flushright}
Qual é o efeito de executar a ação abaixo no início de um programa?

\begin{minted}{haskell}
hSetBuffering stdout NoBuffering
\end{minted}

\begin{enumerate}
  \item Ela impede que o programa escreva na saída padrão.

  \item Ela faz com que a entrada padrão seja lida caractere por caractere.

  \item Ela configura a saída padrão para não manter conteúdo em \emph{buffer}, fazendo com que as saídas sejam enviadas imediatamente.

  \item Ela converte todo texto impresso para maiúsculas.

  \item Ela configura \haskellinline|stdout| para ler dados de um arquivo.


\end{enumerate}

\questionTitle{62}{Leitura de quantidade fixa com n não positivo}
\begin{flushright}
  \footnotesize
  \textit{\textbf{10: Ações de E/S recursivas}}\\
  \textit{c10-lacos-io-03}\\

\end{flushright}
Analise a função:

\begin{haskellcode}
leLista :: Integer -> IO [Integer]
leLista n
  | n <= 0    = return []
  | otherwise = do
      x <- readLn
      xs <- leLista (n - 1)
      return (x:xs)
\end{haskellcode}

Se \haskellinline|leLista (-2)| for executada, o que acontece?

\begin{enumerate}
  \item O programa falha por erro de tipo, pois \haskellinline|Integer| não pode ser comparado com zero.

  \item Nenhum valor é lido; a ação produz a lista vazia por meio de \haskellinline|return []|.

  \item A função devolve \haskellinline|[-2]|, pois o valor inicial de \haskellinline|n| é incluído na lista.

  \item A função entra em recursão infinita, pois \haskellinline|n| nunca se torna igual a zero.

  \item A função tenta ler dois valores e depois devolve uma lista vazia.


\end{enumerate}

\questionTitle{63}{Comentários de linha e de bloco}
\begin{flushright}
  \footnotesize
  \textit{\textbf{3: Expressões e definições}}\\
  \textit{c03-comentarios-valores-01}\\

\end{flushright}
Assinale a alternativa que descreve corretamente a sintaxe de comentários em Haskell.

\begin{enumerate}
  \item Comentários de bloco em Haskell não podem conter outros comentários de bloco em seu interior.

  \item Comentários são avaliados pelo GHCi como expressões do tipo \haskellinline|String|.

  \item Comentários de linha e comentários de bloco só podem aparecer antes da declaração de módulo.

  \item Um comentário iniciado por \haskellinline|--| deve ser fechado com \haskellinline|--| antes do fim da linha.

  \item Um comentário iniciado por \haskellinline|--| vai até o fim da linha, enquanto um comentário delimitado por \haskellinline|{-| e \haskellinline|-\}| pode ocupar várias linhas e pode ser aninhado.


\end{enumerate}

\questionTitle{64}{A onipresença do Prelude}
\begin{flushright}
  \footnotesize
  \textit{\textbf{2: Ambiente de desenvolvimento}}\\
  \textit{c02-ghci-modulos-01}\\

\end{flushright}
No ambiente GHCi, é possível utilizar funções matemáticas como \haskellinline{sqrt} ou operadores como \haskellinline{+} imediatamente após abrir o terminal, sem a necessidade de comandos adicionais. Qual é a justificativa técnica para isso?

\begin{enumerate}
  \item Essas funções não pertencem a nenhum módulo; elas são palavras-chave imutáveis encravadas no analisador léxico (\emph{parser}) do compilador.

  \item Linguagens funcionais não requerem bibliotecas de base, pois todas as operações aritméticas são resolvidas diretamente pelo \emph{hardware} da CPU.

  \item O GHCi carrega automaticamente, por padrão, o módulo especial \texttt{Prelude}, que contém um grande conjunto de funções e tipos básicos.

  \item O GHCi possui um recurso de compilação JIT que adivinha a intenção do programador e cria as funções matemáticas em tempo de execução.

  \item O comando \texttt{ghci} faz uma varredura no disco rígido buscando qualquer arquivo \texttt{.hs} pré-existente e importa todas as funções encontradas silenciosamente.


\end{enumerate}

\questionTitle{65}{Equivalência entre if e guardas}
\begin{flushright}
  \footnotesize
  \textit{\textbf{5: Expressão condicional}}\\
  \textit{c05-if-vs-guardas-03}\\

\end{flushright}
Qual definição com guardas é equivalente à função abaixo?
\begin{haskellcode}
max2 x y = if x > y then x else y
\end{haskellcode}

\begin{enumerate}
  \item \begin{haskellcode}
max2 x y
  | x >= y    = y
  | otherwise = x
\end{haskellcode}

  \item \begin{haskellcode}
max2 x y = if x > y then y else x
\end{haskellcode}

  \item \begin{haskellcode}
max2 x y
  | x < y     = x
  | otherwise = y
\end{haskellcode}

  \item \begin{haskellcode}
max2 x y
  | x > y     = x
  | otherwise = y
\end{haskellcode}

  \item \begin{haskellcode}
max2 x y
  | x > y     = True
  | otherwise = False
\end{haskellcode}


\end{enumerate}

\questionTitle{66}{Rastreamento de recursividade mútua}
\begin{flushright}
  \footnotesize
  \textit{\textbf{9: Funções recursivas}}\\
  \textit{c09-formas-mutua-01}\\

\end{flushright}
Considere:
\begin{haskellcode}
par n | n == 0 = True
      | otherwise = impar (n - 1)

impar n | n == 0 = False
        | otherwise = par (n - 1)
\end{haskellcode}
Qual é o valor de \haskellinline|par 5|?

\begin{enumerate}
  \item \haskellinline|True|

  \item \haskellinline|False|

  \item erro de tipo

  \item não termina

  \item \haskellinline|5|


\end{enumerate}

\questionTitle{67}{Enumeração com passo e consumo finito}
\begin{flushright}
  \footnotesize
  \textit{\textbf{6: Estruturas de dados básicas}}\\
  \textit{c06-listas-operacoes-03}\\

\end{flushright}
Considere a expressão:
\begin{minted}{haskell}
take 3 [2, 2 + 5 ..]
\end{minted}
Qual é o seu resultado?

\begin{enumerate}
  \item \haskellinline{ [2,7,12] }

  \item \haskellinline{ [2,5,10] }

  \item \haskellinline{ [7,12,17] }

  \item A expressão produz a lista infinita completa antes de aplicar \haskellinline{take}.

  \item \haskellinline{ [2,7] }


\end{enumerate}

\questionTitle{68}{Separação entre leitura e soma pura}
\begin{flushright}
  \footnotesize
  \textit{\textbf{10: Ações de E/S recursivas}}\\
  \textit{c10-separacao-puro-01}\\

\end{flushright}
Considere a versão que separa a leitura da soma:

\begin{haskellcode}
main :: IO ()
main = do
  lista <- lerLista
  print (sum lista)

lerLista :: IO [Integer]
lerLista = do
  x <- readLn
  if x == 0
    then return []
    else do
      resto <- lerLista
      return (x:resto)
\end{haskellcode}

Qual é a principal vantagem dessa organização?

\begin{enumerate}
  \item Ela obriga o compilador a transformar \haskellinline|lerLista| em um laço imperativo.

  \item Ela elimina completamente o uso de \haskellinline|IO| do programa.

  \item Ela evita a necessidade de um caso base, pois \haskellinline|sum| detecta automaticamente o sentinela.

  \item Ela faz com que \haskellinline|sum| leia diretamente os números do teclado.

  \item Ela isola a interação com o usuário em \haskellinline|lerLista| e permite que o processamento, como \haskellinline|sum lista|, permaneça puro.


\end{enumerate}

\questionTitle{69}{Inferência de literais no GHCi (parametrizado)}
\begin{flushright}
  \footnotesize
  \textit{\textbf{4: Tipos de Dados}}\\
  \textit{c04-ghci-inferencia-01}\\

\end{flushright}
Qual tipo o GHCi inferiria para o literal \haskellinline{3.14 } ao usar o comando \texttt{:type}?

\begin{enumerate}
  \item Ocorre um erro de sintaxe.

  \item \haskellinline{Num a => a}

  \item \haskellinline{Double}

  \item \haskellinline{Float}

  \item \haskellinline{Fractional a => a}


\end{enumerate}


\end{multicols}
\makeanswersheet{UFOP}{BCC222: Programação Funcional}{José Romildo}{2026/1}{Primeira Prova}{08}{69}{25}{7}
\gabarito{08}{B,C,B,D,B,B,B,C,C,A,E,C,B,B,A,E,E,C,A,A,B,C,E,D,C,D,E,A,B,B,E,E,A,B,B,D,B,C,E,A,C,A,B,C,B,D,C,B,D,D,A,E,A,E,C,C,E,E,E,C,C,B,E,C,D,B,A,E,E}
\end{document}
